% WayFinder Solution Summary -- LaTeX report source.
% Compile with:  pdflatex solution_summary.tex   (or upload to Overleaf)
% A pre-rendered PDF (solution_summary.pdf) ships alongside this file.
\documentclass[11pt,a4paper]{article}

\usepackage[utf8]{inputenc}
\usepackage[T1]{fontenc}
\usepackage[margin=24mm]{geometry}
\usepackage{lmodern}
\usepackage{microtype}
\usepackage{enumitem}
\usepackage{parskip}
\usepackage{titlesec}

\titleformat{\section}{\normalfont\large\bfseries}{\thesection.}{0.6em}{}
\setlist[enumerate]{leftmargin=1.5em,itemsep=4pt,topsep=4pt}

\title{\vspace{-6mm}\textbf{WayFinder: An AI Air Travel Companion}\\[2pt]
  \large\itshape Solution Summary \textperiodcentered\ Problem Statement 1}
\author{Ishika Sattawan}
\date{Expedia Group Campus Hackathon 2026 \textperiodcentered\ Problem Statement 1: AI Air Travel Companion}

\begin{document}

\begin{center}
  {\footnotesize\scshape Expedia Group Campus Hackathon 2026 \textperiodcentered\ Innovation Round}
\end{center}
\vspace{2mm}
\maketitle
\vspace{-4mm}
\hrule
\vspace{4mm}

\begin{abstract}
\noindent WayFinder is a hybrid neuro symbolic travel assistant that reads the traveler, not
just the query. It fuses a traveler's structured profile with signals mined from their messy,
unstructured booking history, searches fifty thousand real itineraries with route level
optimization, and explains every recommendation with a cited receipt. A self grading harness
verifies forty two of forty two expected behaviors across the six provided benchmark prompts,
fully deterministically, with the language model entirely optional.
\end{abstract}

\section{Problem Statement}
Problem Statement 1 of the Innovation Round asks for an AI powered air travel assistant that
suggests and optimizes flight options based on user profiles and inferred travel preferences.
The solution should support multi city travel, direct versus connecting trade offs, flexible
date constraints, and seasonal traffic or pricing patterns, and it should turn all of that into
personalized, optimized, and clearly explained flight recommendations.

\section{The User and Business Problem}
Flight search today treats everyone identically. A solo backpacker chasing the cheapest red eye
and a family of three with a stroller receive the same ranked list. The catch is that a
traveler's real constraints live in two different places. Half of them sit in structured profile
fields, and the other half are buried in how the person actually behaves, in throwaway lines like
``took a 7 hour layover in Singapore to save 120 dollars'' or ``kids melt down at night.'' When a
search engine ignores that second half, it produces generic results, and generic results mean
abandoned searches, avoidable support tickets, and lost bookings. For a platform like Expedia,
personalization that can explain itself is not a nicety. It is directly tied to conversion and to
customer trust.

\section{Proposed Solution}
WayFinder is a hybrid neuro symbolic travel strategist, built from six parts.
\begin{enumerate}
  \item \textbf{Preference Fusion Engine.} It merges the structured profile columns with signals
    mined from messy free text history into provenance tagged preferences. Every inferred
    preference carries the evidence it came from. Contradictions, such as a sixty six year old
    whose history reads ``broke student,'' are detected, resolved by a behavioral evidence wins
    policy, and shown to the user rather than hidden.
  \item \textbf{Route Intelligence.} A deterministic search core runs over fifty thousand
    itineraries. It performs date window search against a simulated travel clock, composes its own
    connections when no published fare serves a request, honors timezone correct weekday patterns
    such as ``Tuesday meeting, back Thursday,'' optimizes visit order for multi city trips through
    permutation and beam search, and discovers open ended regional trips such as a multi city Asia
    loop. When constraints cannot be met, a relaxation ladder loosens them one step at a time and
    narrates every adjustment.
  \item \textbf{Personalized Explainable Ranking.} Scoring weights are derived from the profile
    through a documented mapping, so price sensitivity raises the price weight, a strong direct
    preference raises stop penalties, and a business purpose raises reliability. Every
    recommendation ships with a fit score from zero to one hundred, a per feature breakdown, and
    evidence cited reasons that quote the traveler's own history.
  \item \textbf{Trade off and Market Intelligence.} Every answer surfaces the Recommended,
    Cheapest, and Fastest options with a cost per hour framing, seasonal premiums computed from
    the dataset's own medians, seat scarcity alerts, and shift your dates counterfactuals.
  \item \textbf{Conversational Refinement Loop.} After any plan, follow ups such as ``make it
    cheaper,'' ``no redeyes,'' or ``under 900 dollars'' are parsed into an intent patch and the
    plan re runs, with every applied change surfaced as a chip. Impossible follow ups degrade
    honestly to the closest option, so a refinement never dead ends and never silently forgets an
    earlier constraint.
  \item \textbf{AI Assisted, Deterministically Guarded.} A local language model, auto detected at
    startup, fills natural language gaps and polishes explanations. It is guarded by deterministic
    fallbacks and a number integrity check that rejects any hallucinated figure. With the model
    switched entirely off, the full system, including every benchmark, still passes.
\end{enumerate}

\section{Datasets and Inputs Used}
The solution is built entirely on the provided hackathon data. \texttt{flights\_data.csv}
supplies fifty thousand priced itineraries across thirty five airports and one thousand one
hundred and seventy two routes, each carrying price, seats remaining, on time performance,
season, and demand level. \texttt{user\_data.csv} supplies fifty traveler profiles, each with
sixteen structured columns plus a raw history free text field that holds the messy, contradictory,
real world signal the engine mines. \texttt{benchmark\_prompts.json} supplies six judge prompts,
B01 to B06, each with a list of expected behaviors. Because the flight data is historical,
WayFinder runs on a simulated travel clock set to August 2025, which is how a vague request for
next month becomes a real, bookable window. The system uses no external APIs and runs fully
offline.

\section{Evidence It Works}
A self grading harness executes all six benchmark prompts and programmatically verifies every
expected behavior listed in the benchmark file. Forty two of forty two behaviors pass, fully
deterministically, with each response returning in under half a second, and with thirty eight
unit and end to end tests alongside, including the refinement loop. The prototype itself is a
polished React web application on a FastAPI backend. It presents a traveler dossier with
provenance tagged evidence chips grouped by constraint strength, airline style itinerary cards
with animated fit score rings, an offline world map that draws great circle arcs, a price
calendar, a conversational refine bar, and a live benchmark self grading tab that runs the judges'
own prompts on screen.

\section{Expected Value and Impact}
For travelers, WayFinder delivers recommendations that respect real constraints, such as seats
for the whole family, layover caps, and school holidays, and that explain themselves in plain
language. For the platform, the same capability drives higher search to book conversion through
genuine personalization, fewer support escalations through expectation setting on holiday premiums
and seat scarcity, and an auditable ranking system whose every decision traces to evidence rather
than to a black box. That auditability is precisely what survives regulatory and customer trust
review, which makes the approach not only compelling for a demo but defensible in production.

\end{document}
